\def\stat{ostrikova}





\def\tit{АНАЛИЗ СОВМЕСТНОГО ИСПОЛЬЗОВАНИЯ СТРАТЕГИЙ ЭНЕРГОСБЕРЕЖЕНИЯ 


ДЛЯ~УСТРОЙСТВ~5G С~ОГРАНИЧЕННЫМ ФУНКЦИОНАЛОМ$^*$}





\def\titkol{Анализ совместного использования стратегий энергосбережения 


для~устройств~5G} % с~ограниченным функционалом}





\def\aut{В.\,А.~Бесчастный$^1$, Е.\,С.~Голос$^2$, Д.\,Ю.~Острикова$^3$, 


Е.\,А.~Мачнев$^4$, В.\,С.~Шоргин$^5$, Ю.\,В.~Гайдамака$^6$}





\def\autkol{В.\,А.~Бесчастный, Е.\,С.~Голос, Д.\,Ю.~Острикова и~др.}


%$^3$,  Е.\,А.~Мачнев$^4$, В.\,С.~Шоргин$^5$, Ю.\,В.~Гайдамака$^6$}





\titel{\tit}{\aut}{\autkol}{\titkol}





\index{Бесчастный В.\,А.}


\index{Голос Е.\,С.}


\index{Острикова Д.\,Ю.} 


\index{Мачнев Е.\,А.}


\index{Шоргин В.\,С.}


\index{Гайдамака Ю.\,В.}


\index{Beschastnyi V.\,A.}


\index{Golos E.\,S.}


\index{Ostrikova D.\,Yu.}


\index{Machnev E.\,A.}


\index{Shorgin V.\,S.}


\index{Gaidamaka Yu.\,V.}





{\renewcommand{\thefootnote}{\fnsymbol{footnote}}


\footnotetext[1]


{Исследование выполнено за счет гранта Российского научного фонда №\,23-79-10084 


({\sf https://rscf.ru/project/23-79-10084}).}} 





\renewcommand{\thefootnote}{\arabic{footnote}}


\footnotetext[1]{Российский университет дружбы народов им.\ Патриса Лумумбы, beschastnyy-va@rudn.ru}


\footnotetext[2]{Российский университет дружбы народов им.\ Патриса Лумумбы, 1142220130@rudn.ru}


\footnotetext[3]{Российский университет дружбы народов им.\ Патриса Лумумбы, ostrikova-dyu@rudn.ru}


\footnotetext[4]{Российский университет дружбы народов им.\ Патриса Лумумбы, 1042200071@rudn.ru}


\footnotetext[5]{Федеральный исследовательский центр <<Информатика и~управ\-ле\-ние>> Российской академии наук, 


\mbox{vshorgin@ipiran.ru}}


\footnotetext[6]{Российский университет дружбы народов им.\ Патриса Лумумбы; Федеральный исследовательский 


центр <<Информатика и~управ\-ле\-ние>> Российской академии наук, \mbox{gaydamaka-yuv@rudn.ru}}








 \label{st\stat}





 \thispagestyle{headings}


 


\vspace*{-12pt}





  


  


      


  


  


  \Abst{Недавно стандартизированный тип пользовательского устройства (ПУ) 


с~ограниченным функционалом (Reduced Capability, RedCap) для систем 5G <<новое 


радио>> (New Radio, NR) поддерживает новые механизмы ограничения 


энергопотребления. Используя комбинацию методов сбережения энергии (Discontinuous Reception, DRX), сигнала про\-буж\-де\-ния (Wake-Up Signal, WUS)  


и~RRM-ре\-лак\-са\-ции (Radio Resource Management relaxation), ПУ с~поддержкой 


RedCap может обеспечить высокую энергоэффективность, при этом удовлетворяя 


требованиям к~ско\-рости для большого набора приложений. В~статье исследуется 


влияние перечисленных механизмов на энергопотребление ПУ в~сетях 5G~NR. 


В~качестве основных характеристик рассматриваются энергоэффективность и~время 


полной разрядки батареи. Численные результаты показывают, что совмещение всех 


рассматриваемых механизмов энергосбережения позволяет увеличить время разрядки 


батареи на 15\%--20\%. Механизм WUS более эффективен по сравнению  


с~RRM-ре\-лак\-са\-ци\-ей. При этом указанные выводы сохраняются для широкого 


диапазона плотностей блокаторов и~интенсивностей поступления пакетов.}


  


  \KW{5G; новое радио; DRX; RedCap; WUS; RRM-релаксация}


  


\DOI{10.14357\!/\!08696527230407}{KATMLB}  





\vskip 10pt plus 9pt minus 6pt





 \thispagestyle{headings}


 


 \vspace*{-16pt}


  


\section{Введение}





\vspace*{-4pt}








     Ожидается, что системы 5G <<новое радио>>  


обеспечат большую скорость в~сети доступа для приложений будущего. 


Несмотря на требования, предъявляемые ITU-R в~M.2410~[1], технология 


5G~NR показала достаточно высокую энергоемкость~[2, 3] из-за 


использования широкой полосы пропускания, нескольких уровней MIMO 


(Multiple Input Multiple Output) и~механизма множественного 


подключения. Для устранения этого недостатка консорциумом 3GPP (3rd Generation Partnership Project)


в~релизе~16 было стандартизировано ПУ 


с~ограниченной пропускной способностью 


и~расширены его функциональные возможности в~релизе~17~[4]. При этом 


предполагается, что ПУ RedCap будут использоваться для различных 


приложений, таких как AR/X-VR, автономные камеры видеонаблюдения 


с~детектором движения и~т.\,п.





     


     Для экономии энергии RedCap предлагает несколько механизмов. 


В~част\-ности, 3GPP в~релизе~15 значительно улучшили механизмы DRX, 


используемые в~системах 4G~LTE (4th Generation Long-Term Evolution), что позволило применять их для 


радиоинтерфейса 5G~NR с~узкой диаграммой направленности антенны~[5]. 


Кроме того, в~релизах~16 и~17 были предложены дополнительные 


механизмы, включая сигнал пробуждения 


и~RRM-ре\-лак\-са\-цию~[6]. 


Однако до сих пор исследовательское сообщество не уделяло большого 


внимания их совместному применению.


     


     В статье исследуется влияние совместного использования методов 


энергосбережения 5G~NR RedCap для широкого спектра приложений, 


которые используют ПУ RedCap. С~этой целью используются инструменты 


процессов восстановления, поглощающих цепей Маркова и~стохастической 


геометрии, с~помощью которых сформулирована общая математическая 


модель ис\-сле\-ду\-емой сис\-те\-мы. В~качестве основных метрик 


используются энергоэффективность и~время полной разрядки 


аккумуляторной батареи. Проводится численный анализ исследуемых 


характеристик для различных комбинаций методов энергосбережения.





\vspace*{-4pt}





\section{Системная модель}





\vspace*{-4pt}





     Один из наиболее важных вариантов использования приложений 


AR\!/\!VR\!/\!XR (augmented reality/virtual reality/extended reality)~--- промышленная автоматизация. В связи с~этим в~работе 


рассматривается типовое производственное помещение квадратной формы 


с~длиной стороны~$L$ и~базовой станцией (БС), расположенной в~точке 


$(0,0)$, как показано на рис.~1. Система обслуживает несколько 


пользователей AR\!/\!VR\!/\!XR, которые используют ПУ с~поддержкой RedCap, 


например VR-очки, для получения услуги связи. Чтобы обеспечить 


энергоэффективную работу, ПУ может использовать несколько режимов 


энергосбережения. Далее в~работе рассматривается пользователь, 


находящийся в~точке $(x(t), y(t))$ на плоскости пола здания в~момент 


времени~$t$. Высоты БС и~ПУ постоянны и~считаются равными~$h_A$ 


и~$h_{U}$ соответственно.


     


\begin{figure} %fig1


\vspace*{1pt}


  \begin{center}


 \mbox{%


 \epsfxsize=120mm 


\epsfbox{ost-1.eps}


 }


\end{center}


\vspace*{-12pt}


\Caption{Системная модель }


\vspace*{-9pt}


\end{figure}





    Предполагается, что путь распространения сигнала на линии прямой 


видимости (Line-of-Sight, LoS) может быть заблокирован различным 


специализированным оборудованием, расположенным в~заводских цехах. 


Вероятность блокировки в~рассматриваемом сценарии получена в~[3].


    


    В данной работе учитываются как микро-, так и~макромобильность ПУ. 


Первое относится к~быстрым вращениям ПУ в~руках пользователя~[7, 8], что 


приводит к~потере направленности антенны на БС~[9]. Для учета этого 


эффекта применяется модель, предложенная в~[9], которая параметризуется 


с~использованием статистических данных, представленных в~[10] для 


приложений X--VR. Также используются модели распространения и~антенн, 


описанные в~[3].


     


     В работе рассмотрены четыре комбинации механизмов 


энергосбережения:


     \begin{enumerate}[(1)]


\item DRX. Механизм энергосбережения DRX предполагает работу ПУ 


в~двух циклах: коротком и~длинном. В~начале каждого цикла запускается 


таймер контроля состояния сети. Если данные для передачи поступили на 


ПУ, запускается режим при\-е\-ма\!/\!пе\-ре\-да\-чи данных. Если же 


в~течение короткого цикла данные не были получены, то 


продолжительность следующего цикла увеличивается. Кроме того, если 


в~течение периода мониторинга сети никакие данные не принимаются ПУ, 


в~конце периода ПУ переходит в~спящий режим до следующего цикла. 


Недостаток такого подхода заключается в~том, что циклический мониторинг 


сети также подразумевает энергопотребление, которого невозможно 


избежать;


\item DRX\;+\;WUS. Сигнал WUS представляет собой сигнал пробуждения, 


имеющий синхронизацию с~циклом DRX. В совокупности эти алгоритмы 


работают следующим образом. Пользовательское устройство настроено на работу в~циклах DRX, 


а~также на отслеживание сигнала WUS в~строго заданный период времени. 


Если сигнал WUS не получен в~течение ожидаемого интервала, ПУ 


засыпает и~пропускает цикл DRX. Таким способом можно добиться 


большей экономии энергии за счет сокращения времени мониторинга сети;


\item DRX\;+\;RRM-релаксация. Помимо мониторинга сигнала 


в~нисходящем контрольном канале PDCCH (Packet Downlink Control 


Channel) ПУ также необходимо отслеживать сигналы измерения уровня 


принимаемого сигнала RSRP (Received Signal Reference Power) 


и~RSRQ (Received Signal Reference Quality). В~режиме  


RRM-ре\-лак\-са\-ции наблюдение за этими сигналами отсутствует. 


Ожидается, что эта схема обеспечит хорошую экономию энергии при 


низкой мобильности~ПУ;


\item DRX\;+\;WUS\;+\;RRM-релаксация. Последняя рассматриваемая схема 


объединяет все три механизма сохранения энергии. Эта схема не 


обязательно окажется лучшей, поскольку ожидается, что 


энергоэффективность будет зависеть от мобильности 


и~местоположения~ПУ.


\end{enumerate}





\section{Математическая модель}





    Чтобы оценить энергопотребление ПУ для рассматриваемых схем 


энергосбережения, используется поглощающая цепь Маркова, 


представленная на рис.~2, которая на верхнем уровне имеет четыре 


состояния ($i\hm=1, \ldots , 4$).


    


\begin{figure}[b] %fig2


\vspace*{1pt}


  \begin{center}


 \mbox{%


 \epsfxsize=128mm 


\epsfbox{ost-2.eps}


 }


\end{center}


\vspace*{-11pt}


\Caption{Модель системы в~виде поглощающей цепи Маркова}


\end{figure}





     В предлагаемой модели учтены периоды обслуживания пакетов 


и~периоды простоя. Предполагаем, что эти периоды имеют 


экспоненциальное распределение с~параметрами~$\lambda_D$ (далее 


плотность распределения (ПР) обозначим $f_D(x)$) и~$\lambda_R$ 


с~ПР~$f_R(x)$ соответственно. Потеря активного соединения может 


произойти в~результате микромобильности, макромобильности или перехода в~другой режим RRM. Плотность распределения длительности непрерывного 


соединения ПУ с~БС рассчитывается как минимум из трех случайных 


величин по формуле:


     $$


     f_C(t)=\sum\limits_{a,b,c\in \{M, R, B\}} f_a(t)(1-F_b(t)) (1-F_d(t))\,,


     $$


где $f_M(t)$~--- плотность распределения длительности до наступления 


рассинхронизации лучей, $f_R(t)$~--- до переключения режима RRM, 


а~$f_B(t)$~--- до наступления блокировки прямой видимости. Данные 


распределения предполагаются экспоненциальными 


с~параметрами~$\lambda_M$, $\lambda_R$ и~$\lambda_B$ соответственно.


     


     Вероятности переходов между состояниями верхней цепи условные, 


так как для перехода необходимо, чтобы остальные события произошли 


позже рас\-смат\-ри\-ва\-емого: 


     $$


     w_{i,j}= \int\limits_0^\infty f_j(t) (1-F_y(t)) (1-F_B(t))\,dt\,,\enskip i, j, y 


=1,\ldots , 4\,,\ i\not= j\not= y\,.


     $$


При этом, как можно заметить из рис.~2, нет переходов между состояниями 


мик\-ро\-мо\-биль\-ности, а~ве\-ро\-ят\-ность попадания в~поглощающее состояние 


в~момент~$t$ находится как $1\hm- F_B(t)$. Данный подход воз\-мо\-жен 


благодаря свойству отсутствия последействия. Аналогичным способом 


определяются переходы между состояниями цепи второго уровня.


     


     Мобильность пользователей представлена двумя независимыми 


процессами диффузии по осям~$Ox$ и~$Oy$, ограниченными пределами 


помещения $(0,L)$. Будем считать, что БС находится в~точке $(0, 0)$. 


Статистические свойства пользователя, перемещающегося в~одномерном 


пространстве и~заключенного в~$(0,L)$, описываются функцией Грина 


$p(x,t\vert x_0,t_0)$, представляющей собой вероятность его нахождения 


в~точке~$x$ в~момент~$t$ вместе с~начальными условиями $p(x,t\vert x_0, 


t_0)\hm= \delta(x\hm-x_0)$ и~граничными условиями


     $$


     \fr{\partial p(0,t|x_0,t_0)}{\partial x} = \fr{\partial p(L,t|x_0,t_0)}{\partial x}=0\,.


     $$


     


     Решение задачи может быть получено методом разделения 


переменных, как описано в~[11--13]. В~терминах бесконечного ряда это 


выражение выглядит как


     $$


     p(x_0,t_0) =\fr{1}{L} +\fr{2}{L} \sum\limits_{n=1}^\infty \exp \left[ -\left( 


\fr{n\pi}{2}\right)^2 \fr{t-t_0}{\tau}\right] \cos \left( \fr{nx\pi}{L} \right) 


\cos\left( \fr{nx_0\pi}{L}\right),


     $$


где $\tau=L^2/D$~--- время релаксации; $D$~--- коэффициент диффузии.


     


     На рис.~3 изображена схема перемещения пользователя по нижней 


поверхности помещения из некоторого начального положения $(x_0, y_0)$ 


в~момент времени $t_0\hm=0$ в~положение $(x(t), y(t))$ со смещением 


направленных лучей антенн относительно друг друга.


     





\begin{figure} %fig3


\vspace*{1pt}


  \begin{center}


 \mbox{%


 \epsfxsize=60.916mm 


\epsfbox{ost-3.eps}


 }


\end{center}


\vspace*{-11pt}


\Caption{Смещение луча в~трехмерном пространстве}


\end{figure}





     


     Рассмотрим случайную величину $X(t)$, представляющую 


местоположение пользователя только по оси $Ox$ в~момент времени~$t$. 


Можно заметить, что расстояние~$l_2$ до БС и~угол~$\beta_V$ 


рассогласования антенны в~момент~$t$ можно выразить как


     $$


     l_2(t)= \sqrt{[X(t)]^2+(h_A-h_U)^2}\,;\enskip


     \beta_V= \fr{X(t)}{\sqrt{[x(t)]^2+(h_A-h_U)^2}}\,.


     $$


Обозначая местоположение пользователя по оси~$Y$ случайной величиной 


$Y(t)$, можно заметить, что $\mathrm{tan}\left( 1\hm- \beta_H(t)\right) \hm= 


Y(t)/x_0$. Это приводит к~$$


\beta_H(t) =1-\arctan \left( \fr{Y(t)}{x_0}\right).


$$


Спектральная эффективность в~момент времени~$t$ может быть записана как


$$


S_E(t)=  1+\fr{P_T G_A(t) G_U(t) L(y,t)}{N_0B+I_M}\,,


$$


где $G_A(t)$ и~$G_U(t)$~--- коэффициенты усиления антенн БС и~ПУ 


с~чис\-лом антенных элементов~$N_A$ и~$N_U$ соответственно в~момент 


времени~$t$; $L(y,t)$~--- потери в~момент времени~$t$; $N_0$~--- шум.


     


     Пусть $B(t)$~--- состояние траектории LoS в~момент~$t$, причем~1 


и~0 означают блокированное и~незаблокированное состояния 


соответственно. Чтобы охарактеризовать динамику блокировки, необходимо 


оценить матрицу


     $$


     P(t)= \left( p_{00}(t)\ \  1 - p_{00}(t)\ \ p_{10}(t)\ \ 1- p_{10}(t)\right).


     $$


Используя условную вероятность, получаем


$$


p_{00}= \fr{\mathrm{Pr}\{B(t)=0,\ B(0)=0\}}{\mathrm{Pr}\{B(0)=0\}}\,;\enskip 


p_{10}= \fr{\mathrm{Pr}\{B(t)=0,\ B(0)=1\}}{\mathrm{Pr}\{B(0)=1\}}\,.


$$


Так как процесс блокировки предполагается пуассоновским в~$R^2$, то 


знаменатель можно записать как


$$


\mathrm{Pr}\{B(0)=0\} =e^{-2\lambda_B r_B x_0}.


$$


     


     Используя свойство независимости пуассоновского точечного 


процесса, получаем, что $p_{00}(t)$ полностью определяется 


мобильностью~ПУ:


     $$


     p_{00}(t)= e^{-\lambda_B(2r_B\sqrt{E([X(t)]-x_0)^2+(E[Y(t)])^2-S}}\,,


     $$


где $S=r_B^2 \mathrm{cot}\,(\alpha/2)$~--- площадь пересечения зон блокировки LoS. 





Аналогично можно найти $p_{10}(t)$ как


$$


p_{10}(t)= \fr{\left(1-e^{-\lambda_B (2r_B x_0-S)}\right)e^{-\lambda_BS}}


{1-e^{-2\lambda_B r_B x_0}}\,e^{-\lambda_B \left(2r_B\sqrt{E([X(t)]-


x_0)^2+(E[Y(t)])^2-S}\right)}.


$$


     


     Заметим, что площадь~$S$ совместного пересечения зон блокировки 


LoS зависит от угла~$\alpha$, который сам является функцией от 


времени~$t$ из-за макромобильности ПУ. Используя теорему косинусов, 


определяем искомый угол как


     $$


     \alpha(t)= \arccos \fr{a^2+b^2-c^2}{2ab}\,,


     $$


где параметры определяются как


$$


a=\sqrt{(X(t))^2+(Y(t))^2}\,;\ b=\sqrt{(X(t)-x_0)^2+(Y(t))^2}\,;\ c=\sqrt{x_0^2}\,.


$$





\section{Численный анализ}


     


     В данном разделе проводится численный анализ влияния  


рас\-смат\-ри\-ва\-емых комбинаций механизмов энергосбережения  


с~па\-ра\-мет\-ра\-ми сис\-те\-мы, представленными в~таб\-ли\-це,  


на энер\-го\-эф\-фек\-тив\-ность передачи данных и~время полной разрядки 


аккумуляторной батареи ПУ.





\begin{table}\small


\begin{center}


\tabcolsep=3pt


\begin{tabular}{|c|l|c|}


\multicolumn{3}{c}{Системные параметры}\\


\multicolumn{3}{c}{\ }\\[-6pt]


\hline


\tabcolsep=0pt\begin{tabular}{c}\textbf{Обозна-}\\


\textbf{чение}\end{tabular}&\multicolumn{1}{|c|}{\textbf{Описание}}&\tabcolsep=0pt\begin{tabular}{c}\textbf{Значения}\\


\textbf{по  умолчанию}\end{tabular}\\


\hline


$\lambda_B$ &Плотность блокаторов&0,3~шт./м$^2$\\


$r_B$ &Радиус блокатора&0,5 м\\


$h_B$ &Высота блокатора&3 м\\


$h_U$ &Высота ПУ&1,5 м\\


$h_A$ &Высота тетрагерцевой БС&4 м\\


$P_T$ &Излучаемая мощность на БС&2 Вт\\


$N_A$ &Число конфигураций антенны БС&$16\times16$\\


$N_U$ &Число конфигураций антенны ПУ&$2\times2$\\


$N_O$, $I_M$ &Шум, интерференция&$-84$~дБи, $-60$~дБи\\


$\Delta\varphi$, $\Delta \theta$ &Средняя скорость смещения вокруг осей&10~град\!/\!с\\


$\Delta x$, $\Delta y$ &Средняя скорость смещения от осей&3 см\!/\!с\\


$D$&Коэффициент диффузии&1\\


$L$ &Ширина помещения&50 м\\


$G_A$, $G_U$&Коэффициенты усиления антенн БС и~ПУ&14,58~дБи, 2,57~дБи\\


$\lambda_D$ &Средняя длительность передачи пакета&0,01 с\\


$\lambda_R$ &Средняя длительность между поступлениями пакетов&1\!/\!60~с\\


\hline


\end{tabular}


\end{center}


\end{table}





\begin{figure}[b] %fig4


\vspace*{1pt}


  \begin{center}


 \mbox{%


 \epsfxsize=123.977mm 


\epsfbox{ost-4.eps}


 }


\end{center}


\vspace*{-11pt}


\Caption{Энергоэффективность передачи данных на ПУ в~зависимости от начального 


положения устройства~(\textit{а}) и~от скорости микромобильности~(\textit{б}): \textit{1}~--- DRX; \textit{2}~--- DRX\;+\;WUS; \textit{3}~--- 


DRX\;+\;RRM-релаксация; \textit{4}~--- DRX\;+\;WUS\;+\;RRM-релаксация}


\end{figure}


     


     На рис.~4,\,\textit{а} представлен график зависимости энергоэффективности 


загрузки данных на ПУ от начального положения устройства. Как видно из 


рис.~4,\,\textit{а}, энергоэффективность практически линейно падает с~увеличением 


расстояния от пользователя до БС. При этом каждый из механизмов 


энергосбережения подтверждает свою эффективность: наибольший выигрыш 


дает применение WUS, который позволяет улучшить энергоэффективность 


до~10\%, а совместное применение всех механизмов~--- до~15\%. Обратим 


внимание, что использование механизма WUS более эффективно по 


сравнению с~механизмом RRM-ре\-лак\-са\-ции.


     


     В то же время сильное влияние на энергоэффективность оказывает 


микромобильность пользователя, как показано на рис.~4,\,\textit{б}. Здесь следует 


отметить, что при сохранении того же соотношения в~эффективности 


рассматриваемых механизмов повышение скорости микромобильности 


в~2~раза сокращает энергоэффективность более чем на треть.


     





     


     На рис.~5,\,\textit{а} показана зависимость времени разрядки аккумуляторной 


батареи ПУ от интенсивности поступления пакетов. Более высокая 


интенсивность приводит к~коротким периодам простоя, что влечет за собой 


низкую эффективность стратегий энергосбережения. Время полной разрядки 


батареи сокращается до~20\% при двукратном увеличении нагрузки. При 


этом разница в~эффективности WUS и~RRM-релаксации остается примерно 


на том же уровне.


     


     Наконец, на рис.~5,\,\textit{б} показана зависимость времени разрядки 


аккумулятора ПУ от плотности блокаторов. Увеличение времени разрядки 


также ожидаемо вследствие более частых простоев в~передаче данных, 


возникающих по причине блокировки прямой видимости между ПУ и~БС. 


     


\begin{figure} %fig5


\vspace*{1pt}


  \begin{center}


 \mbox{%


 \epsfxsize=128mm 


\epsfbox{ost-6.eps}


 }


\end{center}


\vspace*{-11pt}


\Caption{Время разрядки аккумуляторной батареи ПУ в~зависимости от интенсивности 


поступления пакетов~(\textit{а}) и~от плот\-ности блокаторов~(\textit{б}): \textit{1}~--- DRX; \textit{2}~--- DRX\;+\;WUS; \textit{3}~--- 


DRX\;+\;RRM-релаксация; \textit{4}~--- DRX\;+\;WUS\;+\;RRM-релаксация}


%\end{figure}


\end{figure}





\section{Заключение}





     В работе предложена математическая модель для сравнения стратегий 


совместного использования методов сбережения энергии, таких как DRX, 


WUS и~RRM-ре\-лак\-са\-ция, в~промышленных условиях для  


VR-при\-ло\-же\-ний. В~качестве основных метрик рассмотрены 


энергоэффективность и~время работы от батареи.


     


     Численные результаты показали, что совмещение всех 


рассматриваемых механизмов энергосбережения наиболее эффективно 


с~точки зрения рассматриваемых характеристик. В~част\-ности, 


использование WUS и~RRM-ре\-лак\-са\-ции позволяет увеличить время 


работы от батареи оконечного устройства на~10\%--15\%. Эти значения 


сохраняются при изменении плотностей блокаторов и~интенсивности 


поступления пакетов. При этом механизм WUS более эффективен по 


сравнению с~RRM-ре\-лак\-са\-цией.


     


     Планируется продолжить исследования с~применением аппарата 


теории ресурсных сис\-тем массового обслуживания~[14] для анализа 


сценариев с~несколькими технологиями радиодоступа при предоставлении 


услуг AR/X--VR.








{\small\frenchspacing


{%\baselineskip=10.8pt


\addcontentsline{toc}{section}{Литература}


\begin{thebibliography}{99}


\bibitem{1-ost}


Minimum requirements related to technical performance for IMT-2020 radio  


interface. ITU-R, 2017. Report  M.2410-0. 


\bibitem{2-ost}


\Au{Beschastnyi V., Ostrikova~D., Moltchanov~D., Gaidamaka~Y., Koucheryavy~Y., 


Samouylov~K.} Balancing latency and energy efficiency in mmWave 5G NR systems with 


multiconnectivity~/\!\!/  IEEE Commun. Lett., 2022. Vol.~26. No.\,8. P.~1952--1956. doi: 


10.1109\!/\!LCOMM.2022.3175663. 


\bibitem{3-ost}


\Au{Ostrikova~D., Beschastnyi V., Moltchanov~D., Gaidamaka~Y., Koucheryavy~Y., 


Samouylov~K.} System-level analysis of energy and performance trade-offs in mmWave 5G NR 


systems~/\!\!/  IEEE T. Wirel. Commun., 2023. Vol.~22. Iss.~11. P.~7304--7318. doi: 


10.1109\!/\!TWC.2023.3250092.


\bibitem{4-ost}


Study on support of reduced capability NR devices (Release 17): Technical Specification 38.875 


V17.0.0.~--- 3GPP, 2021. 136~p. 


{\sf https://\linebreak www.3gpp.org/ftp/Specs/archive/38\_series/38.875/38875-h00.zip}.


\bibitem{5-ost}


Radio Resource Control (RRC) 


protocol specification (Release 15): Technical Specification 38.331 V15.6.0.~--- 3GPP, 2019. 517~p.


{\sf https://www.3gpp.org/ftp/Specs/\linebreak archive/38\_series/38.331/38331-f60.zip}.


\bibitem{6-ost}


\Au{Veedu S., Mozaffari~M., Hoglund~A., Yavuz~E., Tirronen~T., Bergman~J., Wang~J.} 


Toward smaller and lower-cost 5G devices with longer battery life: An overview of 3GPP 


release 17 RedCap~/\!\!/  IEEE Communications Standards Magazine, 2022. Vol.~6. No.\,3. P.~84--90. doi: 


10.1109\!/\!MCOMSTD.0001.2200029.





\bibitem{8-ost} %7


\Au{Petrov V., Moltchanov~D., Koucheryavy~Y, Jornet~J.\,M.} The effect of small-scale 


mobility on terahertz band communications~/\!\!/  5th ACM Conference (International) on 


Nanoscale Computing and Communication Proceedings.~--- New York, NY, USA: ACM, 2018. 


Art.~40. 2~p. doi: 10.1145\!/\!3233188.3242902.





\bibitem{7-ost} %8


\Au{Singh R., Sicker~D.} Parameter modeling for small-scale mobility in indoor THz 


communication~/\!\!/  Global Communications Conference.~--- Piscataway, NJ, USA: IEEE, 


2019. Art. 9013838. 6~p. doi: 10.1109\!/\!GLOBECOM38437.2019.9013838. 





\bibitem{9-ost}


\Au{Petrov V., Moltchanov~D., Koucheryavy~Y., Jornet~J.\,M.} Capacity and outage of 


terahertz communications with user micro-mobility and beam misalignment~/\!\!/  IEEE T. Veh. 


Technol., 2020. Vol.~69. No.\,6. P.~6822--6827. doi: 10.1109\!/\!TVT.2020.2988600.


\bibitem{10-ost}


\Au{Stepanov N.\,V., Moltchanov~D., Begishev~V., Turlikov~A., Koucheryavy~Y.} Statistical 


analysis and modeling of user micromobility for THz cellular communications~/\!\!/  IEEE T. 


Veh. Technol., 2021. Vol.~71. No.\,1. P.~725--738. doi: 10.1109\!/ TVT.2021.3124870.











\bibitem{12-ost} %11


\Au{Bickel T.} A~note on confined diffusion~/\!\!/  Physica~A, 2007. Vol.~377. No.\,1. P.~4--32. doi: 10.1016/j.physa.2006.11.008.





\bibitem{11-ost} %12


\Au{Ezekoye O.\,A.} Conduction of heat in solids~/\!\!/ SFPE handbook of fire protection 


engineering~/ Eds. M.\,J.~Hurley, D.~Gottuk, J.\,R.~Hall, \textit{et al.}~--- 5th ed.~--- New York, NY, USA: Springer, 2016. 


P.~25--52.  


doi: 10.1007\!/\!978-1-4939-2565-0\_2.





\bibitem{13st}


\Au{Mortensen K.\,I., Flyvbjerg~H., Pedersen~J.\,N.}Confined Brownian motion tracked with 


motion blur: Estimating diffusion coefficient and size of confining space~/\!\!/ Frontiers Physics, 


2021. Vol.~8. Art.~583202. 15~p. doi: 10.3389\!/\!fphy.2020.583202.


\bibitem{14st}


\Au{Горбунова А.\,В., Наумов~В.\,А., Гайдамака~Ю.\,В., Самуйлов~К.\,Е.} Ресурсные 


системы массового обслуживания как модели беспроводных сис\-тем связи~/\!\!/  


Информатика и~её применения, 2018. Т.~12. Вып.~3. С.~48--55. doi: 


10.14357\!/ 19922264180307.


   \end{thebibliography}


} }











\vspace*{-6pt}





\hfill{\small\textit{Поступила в~редакцию 15.09.23}}





\vspace*{8pt}





\hrule





\vspace*{2pt}











\hrule





\vspace*{8pt}





%\newpage





%\vspace*{-28pt}





\def\tit{ANALYSIS OF JOINT USAGE OF~ENERGY CONSERVATION STRATEGIES  


FOR~5G DEVICES WITH REDUCED CAPABILITY}











\def\titkol{Analysis of joint usage of~energy conservation strategies  


for~5G devices} %with reduced capability}





\def\aut{V.\,A.~Beschastnyi$^1$, E.\,S.~Golos$^1$, D.\,Yu.~Ostrikova$^1$, E.\,A.~Machnev$^1$,  


V.\,S.~Shorgin$^2$, and~Yu.\,V.~Gaidamaka$^{1,2}$}





\def\autkol{V.\,A.~Beschastnyi, E.\,S.~Golos, D.\,Yu.~Ostrikova, et al.}


%E.\,A.~Machnev$^1$,   V.\,S.~Shorgin$^2$, and~Yu.\,V.~Gaidamaka$^{1,2}$}





\titel{\tit}{\aut}{\autkol}{\titkol}





\vspace*{-8pt}





 \noindent


    $^1$Peoples' Friendship University of Russia (RUDN University), 6~Miklukho-\linebreak


    $\hphantom{^1}$Maklaya Str., Moscow 117198, Russian Federation 





\noindent


$^2$Federal Research Center ``Computer Science and Control'' of the Russian Academy\linebreak


$\hphantom{^1}$of Sciences,   44-2~Vavilov Str., Moscow 119133, Russian Federation








\def\leftfootline{\small{\textbf{\thepage}


\hfill Sistemy i Sredstva Informatiki~---


Systems and Means of Informatics 2023 vol~33 no\,4}


}%


 \def\rightfootline{\small{Sistemy i Sredstva Informatiki~---


 Systems and Means of Informatics 2023 vol~33 no\,4


\hfill \textbf{\thepage}}}





\vspace*{3pt}








  


  \Abste{The recently standardized Reduced capability (RedCap) user equipment (UE) type for 5G  systems,  New Radio (NR), supports new mechanisms to limit


  power consumption. Using a~combination of Discontinuous Reception (DRX), Wake-Up Signal (WUS), 


  and Radio Resource Management (RRM) relaxation  techniques, RedCap-enabled  UEs can provide high energy efficiency while meeting the speed requirements


  for a~large set of applications. The present authors investigate the impact of the above mechanisms


  on the energy consumption of UEs in 5G NR networks. 


  To this aim, the authors propose a~mathematical model that captures the specifics of these mechanisms and environmental and deployment characteristics.


   Energy efficiency and battery life time are considered as the main characteristics. The numerical results demonstrate that combining all


    the considered energy saving mechanisms allows to improve the battery life time by~15\%--20\%. 


   The WUS mechanism shows better performance as compared to the RRM relaxation. These conclusions hold for 


   a~wide range of input system parameters, including packet interarrival time and blockers density.}


  


  \KWE{5G; New Radio; DRX; RedCap; WUS; RRM relaxation}


   





  


\DOI{10.14357\!/\!08696527230407}{KATMLB}  





\vspace*{-10pt}








\Ack





\vspace*{-3pt}





  \noindent


  The study was funded by the Russian Science Foundation, project No.\,23-79-10084 ({\sf https://rscf.ru/project/23-79-10084}). 





%\vspace*{-3pt}





\renewcommand{\bibname}{\protect\rmfamily References}


%\renewcommand{\bibname}{\large\protect\rm References}





{\small\frenchspacing


{ %\baselineskip=12pt


\addcontentsline{toc}{section}{References}


\begin{thebibliography}{99}


  \bibitem{1-ost-1}


ITU-R.  2017.


  Minimum requirements related to technical performance for IMT-2020 radio interface 


Report M.2410-0. Available at: {\sf https://www.itu.int/pub/R-REP-M.2410-2017} (accessed 


October~23, 2023).


  \bibitem{2-ost-1}


  \Aue{Beschastnyi, V., D.~Ostrikova, D.~Moltchanov, Y.~Gaidamaka, Y.~Koucheryavy, and 


K.~Samouylov.} 2022. Balancing latency and energy efficiency in mmWave 5G NR systems 


with multiconnectivity. \textit{IEEE Commun. Lett.} 26(8):1952--1956. doi: 


10.1109\!/\!LCOMM.2022.3175663. 


  \bibitem{3-ost-1}


  \Aue{Ostrikova, D., V.~Beschastnyi, D.~Moltchanov, Y.~Gaidamaka, Y.~Koucheryavy, and 


K.~Samouylov.} 2023. System-level analysis of energy and performance trade-offs in mmWave 


5G NR systems. \textit{IEEE T. Wirel. Commun.}  22(11):7304--7318. doi: 10.1109\!/\!TWC.2023.3250092.


  \bibitem{4-ost-1}


 3GPP. 2021. Study on support of reduced capability NR devices (Release 17): Technical Specification 


38.875 V17.0.0. 136~p. Available at: {\sf 


https://www. 3gpp.org/ftp/Specs/archive/38\_series/38.875/38875-h00.zip} (accessed October~23, 2023).


  \bibitem{5-ost-1}


  3GPP. 2019. 


Radio Resource Control (RRC) protocol 


specification (Release~15): Technical Specification 38.331 V15.6.0. 517~p. Available at: {\sf 


https://www.3gpp.org/ftp/Specs/archive/38\_series/38.331/38331-f60.zip} (accessed October~23, 2023).


  \bibitem{6-ost-1}


  \Aue{Veedu, S., M.~Mozaffari, A.~Hoglund, E.~Yavuz, T.~Tirronen, J.~Bergman, and 


J.~Wang.} 2022. Toward smaller and lower-cost 5G devices with longer battery life: An 


overview of 3GPP release 17~RedCap. \textit{IEEE Communications Standards Magazine} 


6(3):84--90. doi: 10.1109\!/\!MCOMSTD.0001.2200029.


 


  \bibitem{8-ost-1}


  \Aue{Petrov, V., D.~Moltchanov, Y.~Koucheryavy, and J.\,M.~Jornet.} 2018. The effect of 


small-scale mobility on terahertz band communications. \textit{5th ACM Conference 


(International) on Nanoscale Computing and Communication Proceedings}. New York, NY: 


ACM. Art.\ 40. 2~p. doi: 10.1145\!/\!3233188.3242902.





 \bibitem{7-ost-1}


  \Aue{Singh, R., and D.~Sicker.} 2019. Parameter modeling for small-scale mobility in indoor 


THz communication. \textit{Global Communications Conference}. Piscataway, NJ: IEEE.  Art.\ 9013838.


6~p. doi:10.1109\!/\!GLOBECOM38437.2019.9013838. 





  \bibitem{9-ost-1}


  \Aue{Petrov, V., D.~Moltchanov, Y.~Koucheryavy, and J.\,M.~Jornet.} 2020. Capacity and 


outage of terahertz communications with user micro-mobility and beam misalignment. 


\textit{IEEE T. Veh. Technol.} 69(6):6822--6827. doi: 10.1109\!/\!TVT.2020.2988600.


  \bibitem{10-ost-1}


  \Aue{Stepanov, N.\,V., D.~Moltchanov, V.~Begishev, A.~Turlikov, and Y.~Koucheryavy}. 


2021. Statistical analysis and modeling of user micromobility for THz cellular communications. 


\textit{IEEE T. Veh. Technol.} 71(1):725--738. doi: 10.1109\!/\!TVT.2021.3124870.


  


  \bibitem{12-ost-1} %11


  \Aue{Bickel, T.} 2007. A~note on confined diffusion. \textit{Physica~A} 377(1):24--32. doi: 10.1016\!/\!j.physa.2006.11.008.


  


  \bibitem{11-ost-1} %12


  \Aue{Ezekoye, O.\,A.} 2016. Conduction of heat in solids. \textit{SFPE handbook of fire 


protection engineering}. Eds. M.\,J.~Hurley, D.~Gottuk, J.\,R.~Hall, \textit{et al.} New York, NY: Springer. 25--52. doi: 


10.1007\!/\!978-1-4939-2565-0\_2.





  \bibitem{13-ost-1}


  \Aue{Mortensen, K.\,I., H.~Flyvbjerg, and J.\,N.~Pedersen.} 2021. Confined Brownian 


motion tracked with motion blur: Estimating diffusion coefficient and size of confining space. 


\textit{Frontiers Physics} 8:583202. 15~p. doi: 10.3389\!/\!fphy.2020.583202.


  \bibitem{14-ost-1}


  \Aue{Gorbunova, A.\,V., V.\,A.~Naumov, Y.\,V.~Gaidamaka, and K.\,E.~Samouylov.} 


2018. Re\-surs\-nye sis\-te\-my mas\-so\-vo\-go ob\-slu\-zhi\-va\-niya kak mo\-de\-li  


bes\-pro\-vod\-nykh sis\-tem svya\-zi [Resource queuing systems as models of wireless 


communication systems]. \textit{Informatika i~ee Primeneniya~--- Inform. Appl.} 12(3):48--55.  


doi: 10.14357\!/\!19922264180307. 





\end{thebibliography}


} }





\vspace*{-6pt}





\hfill{\small\textit{Received September 15, 2023}} 





\vspace*{-14pt} 


  


\Contr





\vspace*{-4pt}


  


\noindent


\textbf{Beschastnyi Vitalii A.} (b.\ 1992)~--- Candidate of Science (PhD) in physics and mathematics, 


associate professor, Department of Probability Theory and Cyber Security, Peoples' 


Friendship University of Russia (RUDN University), 


6~Miklukho-Maklaya Str., Moscow 117198, Russian Federation; \mbox{beschastnyy-va@rudn.ru}


  


\vspace*{3pt}


  


  \noindent


\textbf{Golos Elizaveta S.} (b.\ 1998)~--- PhD student, Department of Probability Theory and Cyber 


Security, Peoples' 


Friendship University of Russia (RUDN University), 6~Miklukho-Maklaya Str., Moscow 117198, Russian Federation; 


\mbox{1142210130@rudn.ru}


  


  \vspace*{3pt}


  


  \noindent


\textbf{Ostrikova Daria Yu.} (b.\ 1988)~--- Candidate of Science (PhD) in physics and mathematics, 


associate professor, Department of Probability Theory and Cyber Security, Peoples' 


Friendship University of Russia


(RUDN University), 


6~Miklukho-Maklaya Str., Moscow 117198, Russian Federation; \mbox{ostrikova-dyu@rudn.ru}





\vspace*{3pt}


  


  \noindent


\textbf{Machnev Egor A.} (b.\ 1996)~--- PhD student, Department of Probability Theory and Cyber 


Security, Peoples' 


Friendship University of Russia (RUDN University), 6~Miklukho-Maklaya Str., Moscow 117198, Russian Federation; 


\mbox{1042200071@rudn.ru}





\vspace*{3pt}


  


  \noindent


\textbf{Shorgin Vsevolod S.} (b.\ 1978)~--- Candidate of Science (PhD) in technology, senior scientist, 


Federal Research Center ``Computer Science and Control'' of the Russian Academy of Sciences, 


44-2 Vavilov Str., Moscow 119133, Russian Federation; \mbox{vshorgin@ipiran.ru}





\vspace*{3pt}


  


  \noindent


\textbf{Gaidamaka Yuliya V.} (b.\ 1971)~--- Doctor of Science in physics and mathematics, professor, 


Department of Probability Theory and Cyber Security, Peoples' 


Friendship University of Russia (RUDN University), 6~Miklukho-Maklaya 


Str., Moscow 117198, Russian Federation; senior scientist, Federal Research Center ``Computer 


Science and Control'' of the Russian Academy of Sciences, 44-2~Vavilov Str., Moscow 119133, 


Russian Federation; \mbox{gaydamaka-yuv@rudn.ru}








\label{end\stat}





\renewcommand{\bibname}{\protect\rm Литература}    


   


